\documentclass[10pt,a4paper]{amsart}
\usepackage[margin=19mm]{geometry}
\usepackage{amsmath,amssymb,mathtools}
\usepackage[scr=rsfs,cal=euler]{mathalfa}
\usepackage{xfrac}
\usepackage[unicode,hidelinks,bookmarks=false]{hyperref}
\newcommand{\ie}{i.e.\ }
\newcommand{\cf}{cf.\ }
\newcommand{\resp}{respectively\ }
\newcommand{\Subsection}{\subsection}
\providecommand{\nobreakdash}{\nobreak\hbox{-}}
\setlength{\emergencystretch}{2em}
\multlinegap0pt
\usepackage{bm}
\usepackage{amssymb}
\usepackage{tikz}
\usepackage{enumerate}
\usepackage[shortlabels]{enumitem}
\usepackage{csquotes}
\usepackage{float}
\usepackage{xcolor}
\usepackage{mathtools}
\newtheorem{claim}{Claim}
\newtheorem{prop}{Proposition}
\newtheorem{poc}{Poc}
\title{Improvement on the Erd\H{o}s-Kleitman conjecture via the KKL theorem}
\author{Gennian Ge, Jialuo Wang, Zixiang Xu}
\date{Source: arXiv:2603.18948v1 (CC BY 4.0)}
\begin{document}
\maketitle

\noindent\emph{Source and licence.} Improvement on the Erd\H{o}s-Kleitman conjecture via the KKL theorem. Version arXiv:2603.18948v1 (CC BY 4.0), distributed by arXiv under the Creative Commons Attribution 4.0 International licence, http://creativecommons.org/licenses/by/4.0/, as recorded in the arXiv licence metadata for this exact version and shown on the abstract page https://arxiv.org/abs/2603.18948v1. Full licence notice: \texttt{licenses/arxiv-2603.18948-CC-BY-4.0.txt}.\par\smallskip

\noindent\textit{Editorial scope.} Counted unit: the article's prop claim:FinalLinearInd (source0.tex lines 792--794). The article's complete original proof of it is delivered with it (source0.tex lines 795--862, one \texttt{proof} environment of 68 lines). No mathematics is written, repaired or re-derived: the statement and the proof are the article's own lines, in the article's own notation, and no retained line is substituted or re-flowed.  Retained uncounted as the source-local context the proof needs: lines 663--665; lines 719--721; lines 761--763; lines 784--786.  Nothing was omitted from any retained span, so the package carries no omission record.  No retained line contains a citation, so the wrapper carries no bibliography and no bibliography entry.  No retained line contains a figure environment, an image-inclusion command or drawing code, so no image is delivered and the package declares no asset dependency.  Editorial formatting only, in addition: an AMS \texttt{amsart} preamble that restates the article's own declarations and macros used by the retained lines, namely the environments \texttt{claim}, \texttt{prop}, \texttt{poc} on separate counters, together with the standard packages those lines need.  The retained environments print in this document as Claim 1, Proposition 1, Poc 1 in order of appearance, whereas the article prints the same items with the numbering of its own full document.  The complete native source of the article as distributed in the arXiv e-print for this exact version is bundled unchanged as \texttt{sources/196740/source0.tex}.
\begin{claim}\label{claim:orthogonal of distinct w_k}
    For any distinct \(k_{1}<k_{2}\in\{0,1,2,\ldots,s-1\}\), any polynomial \(f\in\mathcal{W}_{k_{1}}\) is orthogonal to any polynomial \(g\in\mathcal{W}_{k_{2}}\).
\end{claim}

\begin{claim}\label{claim:SudakovsBound}
    For each \(G\in\mathcal{G}\) and each \(k\in\{0,1,\ldots,s-1\}\), take arbitrary one polynomial \(w_{k,G}\in\mathcal{W}_{G}\cap \mathcal{W}_{k}\) respectively, then the polynomials in \(\{w_{k,G}\}_{G\in\mathcal{G},k\in \{0,1,\ldots,s-1\}}\) are linearly independent.
\end{claim}

\begin{equation}\label{containment:BjCj}
    B_{j}\subseteq C_{j}^{G}
\end{equation}

\begin{equation}\label{RelationBDBD}
    \bigg(\bigsqcup_{j=1}^{k}D_{j}^{H}\bigg)\cap B_{h}\neq\emptyset
\end{equation}

\begin{prop}\label{claim:FinalLinearInd}
   The polynomials in \(\{p_{G,k,S}:G\in \mathcal{G}_{0},S\subseteq [s-1],|S|=k\}\sqcup \{q_{H,k}: H\in \mathcal{G}\setminus\mathcal{G}_{0}\}\) defined above are linearly independent. 
\end{prop}

\begin{proof}[Proof of Proposition~\ref{claim:FinalLinearInd}]
We first show two auxiliary claims.
\begin{claim}\label{claim:PPPP}
   For any \(k\in [s-1]\), any distinct \(S\) and \(S'\) with \(|S|=|S'|=k\), and any \(G,G'\in\mathcal{G}_{0}\), \(p_{G,k,S}\) is orthogonal to \(p_{G',k,S'}\).
\end{claim}
\begin{poc}
    {Since \(S\neq S'\) and \(|S|=|S'|=k\), there exists \(h\in S\setminus S'\).
By definition, \(p_{G,k,S}(\bm{x})\) contains the factor \(\prod_{j\in C_h^{G}} x_j\), while \(p_{G',k,S'}(\bm{x})\) contains the factor \(\prod_{j\in C_h^{G'}} (1-x_j)\).
By~\eqref{containment:BjCj}, we have \(B_h\subseteq C_h^{G}\cap C_h^{G'}\), then the pointwise product \(p_{G,k,S}(\bm{x})\cdot p_{G',k,S'}(\bm{x})\) contains
\[
\prod_{j\in B_h} x_j(1-x_j),
\]
which vanishes for every \(\bm{x}\in\{0,1\}^n\).
This finishes the proof.}
\end{poc}
\begin{claim}\label{Claim:PQPQ}
    For any \(k\in [s-1]\), any \(G\in\mathcal{G}_{0}\) and any $H\in \mathcal{G}\setminus \mathcal{G}_{0}$, 
     $q_{H,k}$ is orthogonal to $p_{G,k,S}$ if $S\neq [k]$.
\end{claim}
\begin{poc}
    { Assume \(S\neq [k]\) and choose \(u\in [k]\setminus S\).
By definition, \(p_{G,k,S}(\bm{x})\) contains the factor \(\prod_{j\in C_u^{G}} (1-x_j)\),
while \(q_{H,k}(\bm{x})\) contains the factor \(\prod_{j\in \bigsqcup_{i=1}^{k} D_i^{H}} x_j\).
By \eqref{containment:BjCj} we have \(B_u\subseteq C_u^{G}\), and by \eqref{RelationBDBD} we have
\[
\bigg(\bigsqcup_{i=1}^{k} D_i^{H}\bigg)\cap B_u \neq \emptyset .
\]
Hence there exists \(t\in C_u^{G}\cap \big(\bigsqcup_{i=1}^{k} D_i^{H}\big)\).
Consequently, the pointwise product \(p_{G,k,S}(\bm{x})\, q_{H,k}(\bm{x})\) contains the factor
\(x_t(1-x_t)\), which vanishes for every \(\bm{x}\in\{0,1\}^n\).
Thus \(p_{G,k,S}\cdot q_{H,k}= 0\) on \(\{0,1\}^n\). This finishes the proof.}
\end{poc}

Now suppose that
\begin{equation}\label{eq:Assumption}
    \sum_{G\in \mathcal{G}_{0},|S|=k}\alpha_{G,S}\cdot p_{G,k,S}+\sum_{H\in \mathcal{G}\setminus\mathcal{G}_{0}}\beta_{H}\cdot q_{H,k}=0.
\end{equation}
By Claim~\ref{claim:orthogonal of distinct w_k}, Claim~\ref{claim:PPPP} and Claim~\ref{Claim:PQPQ}, we can see
\[\sum_{G\in \mathcal{G}_{0}}\alpha_{G,[k]}\cdot p_{G,k,[k]}+\sum_{H\in \mathcal{G}\setminus\mathcal{G}_{0}}\beta_{H}\cdot q_{H,k}\] is orthogonal to each $p_{G,k,S}$, where $S\neq [k]$. Therefore, we have
\[
\begin{aligned}
0&=\bigg(\sum_{G\in \mathcal{G}_{0}}\alpha_{G,[k]}\cdot p_{G,k,[k]}+\sum_{H\in \mathcal{G}\setminus\mathcal{G}_{0}}\beta_{H}\cdot q_{H,k}\bigg)\cdot\bigg(\sum_{G\in \mathcal{G}_{0},|S|=k}\alpha_{G,S}\cdot p_{G,k,S}+\sum_{H\in \mathcal{G}\setminus\mathcal{G}_{0}}\beta_{H}\cdot q_{H,k}\bigg)\\
&=\bigg(\sum_{G\in \mathcal{G}_{0}}\alpha_{G,[k]}\cdot p_{G,k,[k]}+\sum_{H\in \mathcal{G}\setminus\mathcal{G}_{0}}\beta_{H}\cdot q_{H,k}\bigg)^2.
\end{aligned}
\]
   It then follows that
     \begin{equation}\label{eq1}
         \sum_{G\in \mathcal{G}_{0}}\alpha_{G,[k]}\cdot p_{G,k,[k]}+\sum_{H\in \mathcal{G}\setminus\mathcal{G}_{0}}\beta_{H}\cdot q_{H,k}=0.
     \end{equation}
     Moreover, by~\eqref{eq:Assumption} and~\eqref{eq1}, we have
      \begin{equation}\label{eq2}
         \sum_{G\in \mathcal{G}_{0},S\in \binom{[s-1]}{k}\setminus {[k]}}\alpha_{G,S}\cdot p_{G,k,S}=0.
     \end{equation}
     For any two distinct $S_1,S_2\in \binom{[s-1]}{k}\setminus \{[k]\}$, by Claim~\ref{claim:PPPP}, we have 
     $$\bigg(\sum_{G\in \mathcal{G}_{0}}\alpha_{G,S_1}\cdot p_{G,k,S_1}\bigg)\cdot \bigg(\sum_{G\in \mathcal{G}_{0}}\alpha_{G,S_2}\cdot p_{G,k,S_2}\bigg)=0.$$
     Moreover, by~\eqref{eq2}, we have
     $$0=\bigg(\sum_{G\in \mathcal{G}_{0}}\alpha_{G,S_1}\cdot p_{G,k,S_1}\bigg)\cdot\bigg(\sum_{G\in \mathcal{G}_{0},S\in \binom{[s-1]}{k}\setminus {[k]}}\alpha_{G,S}\cdot p_{G,k,S}\bigg)=\bigg(\sum_{G\in \mathcal{G}_{0}}\alpha_{G,S_1}\cdot p_{G,k,S_1}\bigg)^2.$$
     Hence, for each $S\in \binom{[s-1]}{k}\setminus\{[k]\}$, we have
     \begin{equation}\label{eq3}
         \sum_{G\in \mathcal{G}_{0}}\alpha_{G,S}\cdot p_{G,k,S}=0.
     \end{equation}
      Note that for each $G\in \mathcal{G}$ there is at most one element of $\mathcal{W}_G$ appearing in~\eqref{eq1}, by the linear independence given by Claim~\ref{claim:SudakovsBound}, we conclude that $\alpha_{G,[k]}=0,\beta_{H}=0$ for all $G\in \mathcal{G}_0$ and $H\in \mathcal{G}\setminus\mathcal{G}_{0}$.
     Similarly, since for each $G\in \mathcal{G}_{0}$ there is at most one element of $\mathcal{W}_G$ appearing in~\eqref{eq3}, we have $\alpha_{G,S}=0$ for any $G\in \mathcal{G}_{0}$ and $S\neq [k]$. Then we can see if assumption~\eqref{eq:Assumption} holds, then all of the coefficients are zero. This finishes the proof.



    
\end{proof}

\end{document}
